\chapter{Conclusion}\label{ch:conclusion}

This thesis assesses three forms of government support for South Korea's crop disaster insurance program: premium subsidies, administrative cost subsidies, and state reinsurance. Each form addresses a different constraint on insurance provision. Their assessment therefore focuses on distinct outcomes: access to insurance and the distribution of subsidies for premium support, cost and service quality for administrative support, and risk allocation and fiscal exposure for reinsurance.

The production comparisons provide limited evidence on how insurance availability changes agricultural production. Across seven annual crops, the mean production contrast is 0.020 log points, with an exploratory 95 percent interval of [\textminus 0.277, 0.317]. Its sign changes with the comparison crops and aggregation weights, while divergent pre-pilot trajectories weaken a causal interpretation. The municipal rice estimate is \textminus 0.021 log points, with a 95 percent interval of [\textminus 0.064, 0.022], and varies with the municipalities included. This comparison holds the crop fixed and shows closer observed pre-pilot trajectories, but covers only two post-introduction years and remains vulnerable to differential local shocks. These estimates establish neither a zero production response nor the absence of moral hazard.

The institutional assessment identifies three priorities: evaluate tapered premium support against distribution and coverage objectives, link administrative payments to comparable measures of cost and service quality, and compare current reinsurance arrangements with a catastrophe-focused alternative. These are options to evaluate, rather than designs shown to be optimal by the production estimates. The broader contribution is to distinguish the purposes of government support and the evidence needed to assess each one. Whether support is well designed depends on who receives protection, how efficiently it is delivered, and how losses are shared between private insurers and the public sector.
